\documentclass[11pt]{article}
\usepackage[margin=1in]{geometry}
\usepackage{amsmath,amssymb,bm}
\usepackage[T1]{fontenc}
\usepackage{lmodern}

\title{Cortical Dynamical System Equations (TPModel)}
\author{}
\date{\today}

\begin{document}
\maketitle

The cortical model used in this repository is a stochastic, discrete-time, linear state-space system:
\begin{align}
\mathbf{x}_{k+1} &= \mathbf{A}\mathbf{x}_k + \mathbf{B}\mathbf{u}_k + \mathbf{w}_k, \\
\mathbf{y}_k &= \mathbf{C}\mathbf{x}_k.
\end{align}

Here,
\begin{align}
\mathbf{x}_k \in \mathbb{R}^{n},\quad
\mathbf{u}_k \in \mathbb{R}^{m},\quad
\mathbf{y}_k \in \mathbb{R}^{(10n)^2}.
\end{align}

\section*{State Matrix Construction}
\begin{align}
\mathbf{A} = \mathbf{A}_{\mathrm{base}} + \kappa\,\widetilde{\mathbf{W}}.
\end{align}

\(\mathbf{A}_{\mathrm{base}}\) is block-diagonal with damped rotation blocks
\begin{align}
\mathbf{R}(\theta,\delta)
= \delta
\begin{bmatrix}
\cos\theta & -\sin\theta \\
\sin\theta & \cos\theta
\end{bmatrix}.
\end{align}

Distance-based coupling is row-normalized:
\begin{align}
\widetilde{W}_{ij} &\propto \exp\!\left(-\frac{d_{ij}^2}{2\sigma_c^2}\right),
\quad \widetilde{W}_{ii}=0,
\quad \sum_j \widetilde{W}_{ij}=1.
\end{align}

\section*{Process Noise}
Spatial covariance:
\begin{align}
Q_{ij} = \sigma_w^2\exp\!\left(-\frac{d_{ij}^2}{2\sigma_Q^2}\right),
\end{align}
with a small diagonal regularization term added in code for positive definiteness.

Temporal coloring uses AR(1):
\begin{align}
\mathbf{z}_k &= \alpha\mathbf{z}_{k-1} + \sqrt{1-\alpha^2}\,\boldsymbol{\varepsilon}_k,
\quad \boldsymbol{\varepsilon}_k\sim\mathcal{N}(\mathbf{0},\mathbf{I}), \\
\mathbf{w}_k &= \mathbf{L}\mathbf{z}_k,
\quad \mathbf{L}\mathbf{L}^\top = \mathbf{Q}.
\end{align}

\section*{Image Observation Model}
Each state component modulates a Gaussian spatial blob in the cortical image:
\begin{align}
C_{p,i} = \exp\!\left(-\frac{(x_p-x_i)^2+(y_p-y_i)^2}{2\sigma_b^2}\right),
\end{align}
where pixel index \(p\) corresponds to a location \((x_p,y_p)\), and \((x_i,y_i)\) is the center of cortical area \(i\).

\section*{Input Matrix Construction}
For channel \(j\) targeting area \(t_j\), input spread is Gaussian in cortical distance:
\begin{align}
B_{ij} &\leftarrow 0.05\,\exp\!\left(-\frac{d_{i,t_j}^2}{2\sigma_u^2}\right) \text{ (normalized)}, \\
B_{t_j,j} &\leftarrow 0.30.
\end{align}

\end{document}
